\scanpage{3}
8. \uline{Une \deleted{\unsure{esti}}estimation} (\uline{heuristique}) \uline{du nombre des valeurs propres de frobenius dans la cohomologie évanescente qui sont des unités $p$-adiques.}

8.1. Soit $X$ un schéma lisse et propre sur le corps parfait $k$ \inserted{de car. $p\handwritten{>}0$}, et soit $W$ l'anneau des vecteurs de Witt sur $k$. Si on dispose d'une bonne théorie \unsure{de} la cohomologie cristalline $H(X)$ (au sens des catégories dérivées), on devra avoir la relation suivante avec la cohomologie de De Rham
\[
\handwritten{H_{\rm DR}}(X)=\underset{\handwritten{\smile}}{\overset{\handwritten{\frown}}{H\,(}}X)\overset{\rm L}{\otimes}_W k
\tag*{(8.1.1)}
\]
(isomorphisme dans la catégorie dérivée des $k$-modules). On en déduit \unsure{une} formule des coéfficients universels: suites exactes
\[
0\handwritten{\longrightarrow} H^i(X)\otimes_W k\handwritten{\longrightarrow} H^i_{\rm DR}(X)\handwritten{\longrightarrow}\mathrm{Tor}^W_1(H^{i+1}(X),k)\handwritten{\longrightarrow}0\;.
\tag*{(8.1.2)}
\]

Ce sont des suites exactes fonctorielles en $X$, tout comme l'isomorphisme (8.1.1). Sauf erreur, \deleted{\illegible{1}} ceci marche dès maintenant pour $p\handwritten{\ne}2$, grâce à Berthelot.
\inserted{\handwritten{[ On conclut, grâce à 8.4.3,}}
\[
\inserted{\handwritten{0\longrightarrow(H^i(X)\otimes_W k)_{ss}\longrightarrow
\underset{\text{(nul si }i>n\text{)}}{H^i(X,\underline O_X)_{ss}}\longrightarrow
\underset{\text{nul si }i\geq n}{({}_p\mathrm{Tors}^{i+1}(\deleted{\unsure{k}}X))_{ss}}\longrightarrow0}}
\]
% margin note written sideways
\margin{\handwritten{et $0\longrightarrow(H^{\unsure{i+n|i}}(X)\otimes_W k)^{ss}\longrightarrow H^{\unsure{i|i-n}}(X,\underline\Omega_X^n)^{ss}\longrightarrow({}_p\mathrm{Tors}^{\unsure{i+n+1|i-n+1}}(X))^{ss}\longrightarrow0$ où l'indice $ss$ signifie la partie semi-simple pour $F_X=F_X^*\sigma$ et \deleted{où} l'exposant $ss$ \rule[.5ex]{2em}{.3pt} $V_X=\sigma^{-1}F_{X*}$ ]}}

Si maintenant $Y$ est une section hyperplane de \deleted{degré assez élevé de} $X$ (supposée maintenant $\handwritten{\subset}\mathrm P^r_k$), on trouve une suite exacte analogue pour $Y$, et un homomorphisme \inserted{\handwritten{$u=(u_0,u_1,u_2)$}} de la première dans la seconde. Le diagramme du serpent nous donne alors une suite exacte
\[
\inserted{\handwritten{0\longrightarrow\mathrm{Ker}\,u_0\longrightarrow\mathrm{Ker}\,u_1\longrightarrow}}
\handwritten{\mathrm{Ker}\,u_2\longrightarrow}E(Y)\otimes_W k\handwritten{\longrightarrow}
\underset{\handwritten{\substack{\Vert\\\mathrm{Coker}\,u_1}}}{E_{\rm DR}(Y)}
\handwritten{\longrightarrow}\mathrm{Coker}\,u_{\handwritten{2}}\handwritten{\longrightarrow}0\;,
\tag*{(8.1.3)}
\]
où on a posé
\[
E(Y)=\mathrm{Coker}(H^{n-1}(X)\xrightarrow{\handwritten{u_0}}H^{n-1}(Y))\;,\quad
E_{\rm DR}(Y)=\mathrm{Coker}(H_{\rm DR}^{n-1}(X)\xrightarrow{\handwritten{u_1}}H_{\rm DR}^{n-1}(Y))\;,
\]
$X$ étant supposée connexe de dimension $n$, enfin $u_{\handwritten{2}}$ désigne l'application
\[
u_{\handwritten{2}}\mathbin{:}\mathrm{Tor}^{\handwritten{W}}_1(H^n(X),k)\handwritten{\longrightarrow}\mathrm{Tor}^W_1(H^n(Y),k)\;.
\tag*{(8.1.4)}
\]

\deleted{\handwritten{8.2.} Noter que $\mathrm{Tor}^{\deleted{\unsure{n}}W}_1(-,k)$ est isomorphe à $\mathrm{Hom}\deleted{\unsure{(}}_W(k,-)$ (élements de $-$ annulés par $pW$), et que $\deleted{\unsure{H}^{n}(}\quad\mathrm{Tors}^{n}(Y)$ ($=$\deleted{\unsure{(}}sous-module de torsion de $Y$) est dual (si on admet la \uline{dualité de Poincaré} habituelle pour la cohomologie cristalline, y compris les formules pour la torsion) de $\mathrm{Tors}^{n-1}(Y)$, et}
% slanted
\margin{\handwritten{il est \unsure{maladroit|malcommode} d'utiliser ici la dualité de Poincaré \deleted{\illegible{1}}}}
% archive sheet number: 2
